\section{Shell directors and prescribed normals}
\label{sec:shell-directors}

Shell elements require a director field that defines the local through-thickness direction. If no explicit director field is supplied, FEMaster constructs shell normals from the element geometry and equalises neighbouring directions where required by the shell-normal build procedure.

For imported or deliberately oriented shell meshes, the director can be prescribed with an element-nodal field. The field is defined first and then selected with \cmd{NORMAL}.

\begin{syntax}
*FIELD, NAME=<name>, TYPE=ELEMENT_NODAL, COLS=3, FILL=NAN
<element-id>, <local-node>, <nx>, <ny>, <nz>

*NORMAL, FIELD=<name>
\end{syntax}

The local node index is zero-based. The selected field must use the \val{ELEMENT_NODAL} domain and contain exactly three components. \cmd{NORMAL} does not create, normalize, or average the supplied vectors; it only selects the field that is used during shell-normal construction. Using \key{FILL=NAN} is useful for partially prescribed director fields because entries that are not supplied remain distinguishable from explicitly prescribed values and can be completed by the geometric shell-normal construction.

\begin{exampledeck}
*FIELD, NAME=DIRECTORS, TYPE=ELEMENT_NODAL, COLS=3, FILL=NAN
10, 0, 0.0, 0.0, 1.0
10, 1, 0.0, 0.0, 1.0
10, 2, 0.0, 0.0, 1.0
10, 3, 0.0, 0.0, 1.0
*NORMAL, FIELD=DIRECTORS
\end{exampledeck}

The field must be defined before \cmd{NORMAL} in the input deck. The parser resolves enumeration-dependent fields after element topology has been built and before the final model data pass, so the selected director field is available when shell normals are completed.
